\documentclass[11pt,a4paper]{article}
\usepackage{amsmath}
\usepackage{amsfonts}
\usepackage{amssymb}

\title{Gravitational Fields Summary}
\author{Ibrahim}
\date{03/11/2026}

\begin{document}

\maketitle

\section{Formulae}
\subsection{Standard}
\begin{enumerate}

    \item Gravitational Field Strength
    \begin{equation}
        g = \frac{F}{m}
    \end{equation}
    or more useful
    \begin{equation}
        g = -\frac{GM}{r^2}
    \end{equation}
    
    \item Newton's Law of Gravitation
    \begin{equation}
        F_g = \frac{GMm}{r^2}
    \end{equation}

    \item Gravitational Potential
    \begin{equation}
        V_g = -\frac{GM}{r}
    \end{equation}

    \item Change in gravitational potential
    \begin{equation}
        \Delta V_g = \frac{\Delta W}{m}
    \end{equation}

    \item Potential Gradient
    \begin{equation}
        g = -\frac{\Delta V_g}{\Delta r}
    \end{equation}
\end{enumerate}
\clearpage

\subsection{Derived}
\begin{enumerate}
    \item Orbital Time Period
    
    On a planet, $\sum F = F_g$. Hence:

    \begin{equation}
    \begin{aligned}
        F_g &= F_c\\
        \frac{GMm}{r^2} &= m\omega^2r\\
        \omega^2 &= \frac{GM}{r^3}\\
        \frac{4\pi^2}{T^2} &= \frac{GM}{r^3}\\
        T^2 &= \frac{4\pi^2}{GM}r^3\\
\text{Kepler's 3rd Law: } & T^2 \propto r^3\\
        T_{ORB} &= \sqrt{\frac{4\pi^2r^3}{GM}}\\
    \end{aligned}
    \end{equation}

    \item Orbital Velocity ($\neq$ Escape Velocity)

    Using the fact that $m\omega^2r = \frac{mv^2}{r}$
    
    \begin{equation}
    \begin{aligned}
        F_g &= F_c\\
        \frac{GMm}{r^2} &= \frac{mv^2}{r}\\
        \frac{GM}{r} &= v^2\\
        \therefore v_{ORB} &= \sqrt{\frac{GM}{r}}
    \end{aligned}
    \end{equation}
    \clearpage
   
    \item Escape Velocity.
    
    When work is done on a mass to move it to infinity,
    
    the mass' Gain in $E_g =$ Loss in $E_k$. Hence:
    \begin{equation}
    \begin{aligned}
        \Delta W &= \Delta E_k\\
        mV_g &= \frac{1}{2}mv^2\\
        \frac{GM}{r} &= \frac{1}{2}v^2\\
        \therefore v_{ESC} &= \sqrt{\frac{2GM}{r}}\\
        &= \sqrt{2}v_{ORB}
    \end{aligned}
    \end{equation}

    \item Total Energy Change

    \begin{equation}
    \begin{aligned}
        \Delta E_T &= \Delta E_g + \Delta E_k\\
        &= mV_g + \frac{1}{2}mv^2\\
        &= -\frac{GMm}{r} + \frac{GMm}{2r}\\
        &= -\frac{GMm}{2r}\\
        &= -E_k   
    \end{aligned}
    \end{equation}
\end{enumerate}
\clearpage

\section{Nuance(s)}
Be careful with $\Delta$'s.

Eg: Calculating the new radius for a change in gravitational potential

\begin{equation}
\begin{aligned}
    \Delta V_g &= -\frac{GM}{\Delta r}\\
    V_2 - V_1 &\neq -\frac{GM}{r_2-r_1}\\
    &= -GM(\frac{1}{r_2}-\frac{1}{r_1})
\end{aligned}
\end{equation}

Remember $r =$ radius of planet $+$ height above surface
\end{document}